% Generated from the full note by format_quantitative_tex.py; no proof paragraphs omitted.
\section{Explicit regularity and an effective detector from the CM
construction}\label{explicit-regularity-and-an-effective-detector-from-the-cm-construction}

Author/checker: GPT-6 Astra Pro, inherited effort not exposed,
HODGE\_DIVISION /root/sixfolds. Date: 2026-11-21.

This note develops a positive-rank variant of D-SR-01--05. The
regularity argument was suggested by the supplied HECKE\_DIV M7
calculation, but that statement, Markman's output-sheaf theorem, and
Markman's generic-position hypotheses are not premises. All geometry and
regularity needed here are proved below. The deformation step uses the
authorized DEG\_L S-14 input, specifically Perry arXiv:2604.00511v2,
Theorem 1.2 and its proof. Perry remains a preprint in the source-status
convention of this project.

The finite construction below is explicit but has not been executed for
its astronomical torsion order. Its termination follows from the
displayed counting inequalities. Numerical certificates in the
accompanying script concern the symbolic character, matrices, regularity
inequalities, and final integer formula, not an enumeration of the
torsion group.

\subsection{SF-QC01. A strengthened torsion
orbit}\label{SF-QC01}

Let \(C\) be the Fermat plane quartic, \(J=J(C)\), and \(\Theta\) its
principal polarization. Let \(E_i\) be the Gaussian CM elliptic curve.
The three elliptic quotient maps give an isogeny \(j:E_i^3\to J\) with
\(j^*\Theta=2\theta_{\mathrm{prod}}\) and degree eight, as proved in
D-SR-04. Fix a positive squarefree integer \(m\). Put

\[
 P=2^{2048+\lceil\log_2(m+1)\rceil},\quad
 L=\operatorname{lcm}(1,\ldots,P),\quad
 N=4mL^2+1,\quad d=N-1=4mL^2.
\]

For the nine class-number-one radicands \(m\le 163\) one can instead use
the common \(P=2^{2048}\), as checked in D-SR-01. Every prime divisor of
\(N\) exceeds \(P\). In particular \(N\) is odd, and \(j\) induces an
isomorphism \(E_i^3[N]\to J[N]\). There exists a cyclic subgroup
\(H\subset J[N]\) of order \(N\) satisfying both:

\begin{enumerate}
\def\labelenumi{\arabic{enumi}.}
\item
  The translates \(C+h\), \(h\in H\), are pairwise disjoint.
\item
  For every \(P_{0}\in \operatorname{Pic}^{0}(J)\), \[
  H^0(J,I_{Z_H}(2\Theta)\otimes P_0)=0,
  \qquad Z_H=\coprod_{h\in H}(C+h).
  \tag{QC1.1}
  \]
\end{enumerate}

\textbf{Proof.} Sample uniformly from exact-order-N vectors
\(v'\in E_i^3[N]\), and let \(v=j(v')\). D-SR-01 proves that failure of
the first condition has probability less than 216/499. Its proof uses
the bound \(\#((C-C)\cap J[q])\le108q^4\), the lower bound
\(J_{6}(q)>q^{6}/2\) for exact-order-q points, and
\(\sum_{p\mid N}1/p<1/500\). These estimates apply unchanged to the
present m-dependent \(P\).

Identify \(\operatorname{NS}(E_i^{3})\) with integral Hermitian
three-by-three Gaussian matrices using the product principal
polarization. Consider all nonzero such matrices \(M\) with
\(0\le M\le 4I\). There are at most

\[
 B=5^3\,81^3=66\,430\,125
\]

matrices: each diagonal entry is one of 0,1,2,3,4; each off-diagonal
Gaussian entry has real and imaginary parts between \(-4\) and 4.
Positive semidefiniteness actually gives its norm at most 16, so every
nonzero entry is a unit modulo \(N\), because all prime divisors of
\(N\) exceed 16. A row containing such an entry defines a surjective map
\((Z[i]/N)^3\to Z[i]/N\). Consequently \(\operatorname{ker}(M)\)
contains at most \(N^{4}\) points of \(E_i^{3}[N]\), and the probability
that an exact-order-N \(v'\) lies in it is less than \(2/N^{2}\). The
probability of lying in any of these kernels is therefore at most
\(2B/N^{2}<1/4\). The combined failure probability is below
\(216/499+1/4<1\). Choose \(v'\) outside both bad sets, and set
\(H=\langle j(v')\rangle\).

Suppose that (QC1.1) fails. Then an effective divisor \(D\) numerically
equivalent to \(2\Theta\) contains every curve \(C+h\). Write its
distinct irreducible components as \(D_{1},\ldots ,D_k\). Since
\(\Theta ^{2}D=12\) and \(\Theta ^{2}D_i\) is a positive integer,
\(k\le 12\). For every \(h\in H\), both \(D\) and \(D+h\) contain all
the \(N\) curves. If these divisors had no common component, their
proper intersection would be an effective curve of \(\Theta\)-degree 24.
Each curve \(C+h\) has \(\Theta\)-degree three, so \(3N\le 24\), a
contradiction. Thus every \(h\) belongs to at least one set

\[
 S_{ij}=\{h\in H:D_i+h=D_j\}.
\]

Each nonempty \(S_{ij}\) is a coset of
\(H\cap \operatorname{Stab}(D_i)\). If all these subgroups were proper,
each coset would have at most \(N/p_{\min}\) points, where
\(p_{\min}>P\) is the smallest prime divisor of \(N\). The at most
\(k^{2}\le 144\) cosets could not cover \(H\). Therefore some \(D_i\) is
H-stable.

Effective divisors on an abelian variety are nef: for a curve, translate
the divisor to meet it properly and compute the nonnegative
intersection. Hence \(D_i\) and \(2\Theta -D_i\) are nef. The matrix
\(M\) corresponding to \(j^*D_i\) is a nonzero integral Gaussian
Hermitian matrix with \(0\le M\le 4I\). Since \(D_i\) is invariant under
\(v=j(v')\), its polarization homomorphism kills \(v\). Pullback gives
\(Mv'=0\), contrary to the choice of \(v'\). This proves (QC1.1).

The construction is finite: enumerate exact-order-N torsion vectors,
test avoidance of \(C-C\) and the finite displayed matrix kernels, and
take the first successful vector. No test over the continuum
\(\operatorname{Pic}^{0}(J)\) is needed, because the divisor argument
proves its vanishing consequence uniformly.

\subsection{SF-QC02. The coherent output and its complete
character}\label{SF-QC02}

Let \(A=J_x\times \widehat{J}_\alpha\), identify \(\widehat{J}\) with
\(J\) by \(\Theta\) when writing matrices, and put

\[
 x=\Theta_x,\qquad y=\Theta_\alpha,\qquad
 \gamma=c_1(\mathcal P_{x,\alpha}),\qquad
 H_d=dx+y,\qquad H_b=3(x+y),\qquad h=3H_d.
\]

Both \(H_b\) and \(h\) are very ample: the third power of an ample line
bundle on an abelian variety is very ample. Let \(Y=J_y\times A\) and
\(\pi :Y\to A\). Write \(a(y,x)=y+x\). On \(Y\) set

\[
 \mathcal L=\mathcal O(\Theta_y)\otimes a^*\mathcal O(\Theta)
                       \otimes\mathcal P_{y,\alpha}^{-1},
 \qquad D_1=\{y\in Z_H\},\quad D_2=\{y+x\in-Z_H\}.
\]

Define

\[
 \mathcal M=\mathcal L\otimes I_{D_1}I_{D_2},\qquad
 \mathcal G=R\pi_*\mathcal M,\qquad
 E=\mathcal G^\vee[-1].
 \tag{QC2.1}
\]

Then \(E\) is a coherent sheaf of rank \(r=8d\), locally of projective
dimension at most one. It is strongly equivariantly semiregular for a
coherent action of a finite subgroup of \(A[N]\times \widehat{A}[N]\).

To prove the sheaf assertion, the change of variables
\((y,x)\mapsto (y,z=y+x)\) makes \(D_{1},D_{2}\) pullbacks from
independent factors. Their ideals are Tor-independent and their
intersection equals their product. There is therefore an exact sequence
on the total space

\[
 0\longrightarrow\mathcal M\longrightarrow\mathcal L
 \longrightarrow\mathcal L|_{D_1}\oplus\mathcal L|_{D_2}
 \longrightarrow\mathcal L|_T\longrightarrow0,
 \quad T=D_1\cap D_2.
 \tag{QC2.2}
\]

There are \(N^{2}\) disjoint components of \(T\), each isomorphic to
\(C\times C\times \widehat{J}\). Its map \(f:T\to A\) is finite: after
translations and signs, its first-factor map is the addition
\(C\times C\to C^{(2)}\simeq\Theta\subset J\). For a nonhyperelliptic
genus-three curve the Abel map \(C^{(2)}\to\Theta\) is an embedding,
since an effective divisor of degree two has a unique section, and the
differential criterion gives an immersion. The quotient
\(C\times C\to C^{(2)}\) is finite flat of degree two. Thus \(\Theta\)
is smooth and each component of \(f\) factors as a finite flat
degree-two cover of a smooth divisor in A.

Set

\[
 W=\pi_*\mathcal L,\quad
 U=\pi_*(\mathcal L|_{D_1}\oplus\mathcal L|_{D_2}),\quad
 \mathcal T=f_*(\mathcal L|_T).
\]

On every \(\pi\)-fibre, \(L\) has class \(2\Theta\) and no higher
cohomology; \(W\) is locally free of rank eight. Each curve restriction
has degree six, and genus three implies \(H^{1}=0\) and \(h^{0}=4\).
Hence \(U\) is locally free of rank 8N. Since \(f\) is finite, \(T\) has
no higher direct image. Sequence (QC2.2) represents \(G\) by
\([W\to U\to T]\) in degrees 0,1,2. By (QC1.1), \(W\to U\) is fibrewise
injective, so \(Q=U/W\) is a vector bundle of rank \(8N-8=r\) and
\(G=[Q\to T]\) in degrees 1,2.

Each summand of \(T\) is the pushforward of a vector bundle on a smooth
divisor; hence it has local projective dimension one and only its first
sheaf Ext into \(O_A\) is nonzero. Finite duality identifies

\[
 \mathcal T^D:=\mathcal E xt_A^1(\mathcal T,\mathcal O_A)
 =\bigoplus_{i,j} f_{ij,*}
       ((\mathcal L|_{T_{ij}})^{-1}\otimes\omega_{T_{ij}}).
\]

Dualizing \([Q\to T]\) gives a coherent sheaf \(E\) and an exact
sequence

\[
 0\longrightarrow Q^*\longrightarrow E\longrightarrow\mathcal T^D
 \longrightarrow0.
 \tag{QC2.3}
\]

This also proves \(\operatorname{pd}(E)\le 1\). No statement about a
codimension-four nonfree locus is used.

For semiregularity, let \(q:J\to Z=J/H\) and \(F=I_{q(C)}\). D-SR-02
proves injectivity of the trace in degrees zero, one and two for \(F\)
and its transforms by inversion. Thus \(F\boxtimes (-1)^*F\) is fully
semiregular by the three Hochschild Künneth summands. Pullback under
\(q\times q\) gives \(I_{Z_H}\boxtimes I_{-Z_H}\) with its honest
\(H\times H\) descent action. Its invariant \(\operatorname{Ext}^{2}\)
is the original \(\operatorname{Ext}^{2}\), and trace pullback is
injective. Tensor by \(O(\Theta )\boxtimes O(\Theta )\), make the linear
change of variables, apply Fourier--Mukai on one factor, and finally
derived-dual and shift as in (QC2.1). These equivalences or
anti-equivalences preserve trace injectivity and transport the whole
coherent action. Translation and degree-zero tensor functors are taken
to functors of the same kind under these operations. Their action on
isomorphism classes is an automorphism of the abelian variety of such
functors, so the resulting group remains an injective finite subgroup of
\(A\times \widehat{A}\) of exponent \(N\). Coherence is transported from
the honest original descent action, not inferred from invariance of a
K-class. Since the output is a coherent sheaf, its ordinary negative
self-Ext groups vanish.

Here is the full character, independently calculated by exterior
integration. Put \(u^{2}=-d\), \(a=(1+1/u)/2\) and \(b=(1-1/u)/2\). Then

\[
 \operatorname{ch}(I_{Z_H}(\Theta))
 =1+\Theta-\frac d2\Theta^2-\frac d6\Theta^3
 =a e^{u\Theta}+b e^{-u\Theta}.
\]

In Darboux coordinates define

\[
 Q_u=\prod_{j=1}^3(u a_j-U_j)(u b_j-V_j),\qquad
 t=-ab(Q_u+Q_{-u}).
\]

With the sign convention of (QC2.1),

\[
 e^{\gamma/2}\operatorname{ch}(E)=t+P_d(H_d),
 \tag{QC2.4}
\]

where

\[
 P_d(z)=8d+2(d-1)z-z^2-\frac{d-1}{12d}z^3
     +\frac1{48d}z^4+\frac{d-1}{960d^2}z^5
     -\frac1{5760d^2}z^6.
 \tag{QC2.5}
\]

For completeness, in one symplectic plane, integrating the middle
coordinates c,e gives

\[
 \int_{c,e}\exp(s(a+c)(b+e)+tce-cV+eU)
  =(s+t)+st\,ab+s(aV-bU)-UV-s\,abUV.
\]

For opposite roots \(s=-t\) this is \(-e^{\gamma _j/2}(sa-U)(sb-V)\),
producing the two Weil determinant terms. For equal roots \(s=t\) it
equals

\[
 2s\exp\left(\frac{\gamma_j}{2}+\frac{s}{2d}(d a b+UV)\right).
\]

Multiplying the three planes, then dualizing and multiplying by \(-1\)
for the shift, gives the unequal-root term \(-ab(Q_u+Q_{-u})\) and the
equal-root polynomial

\[
 -8u^3\left(a^2e^{-uH_d/(2d)}-b^2e^{uH_d/(2d)}\right).
\]

Expanding through degree six proves (QC2.5). There is no additional
volume factor because \(\Theta\) on \(J\) is principal.

The Weil operator on \(H_{1}(J)\oplus H_{1}(J)\) is

\[
 J_K=\begin{pmatrix}0&-1\\d&0\end{pmatrix},\qquad J_K^2=-d.
\]

Its two eigenspaces on \(H^{1}\) have generators \(ua_j-U_j\),
\(ub_j-V_j\) and their conjugates, each of Hodge signature (3,3). It is
Rosati-skew for the polarization matrix \(\operatorname{diag}(d,1)\);
that product is
\(\bigl(\begin{smallmatrix}0&-d\\d&0\end{smallmatrix}\bigr)\). If
\(L_{0}\subset H_{1}(J,Q)\) is a rational symplectic Lagrangian,
\(L_{0}\oplus L_{0}\) is \(J_K\)-stable and totally isotropic for the
associated K-Hermitian form. Thus the Hermitian space is split with
determinant class \([-1]\). The two determinant terms \(Q_u,Q_{-u}\) are
independent and \(ab=N/(4d)\) is nonzero. Therefore \(t\) is a nonzero
rational primitive Weil class. In particular \(H_d t=0\).

Writing

\[
 a_d=\frac{d-1}{4d},\qquad
 D=\det E,\qquad c_1(D)=2(d-1)H_d-4d\gamma,
\]

the normalized entire character is

\[
 \kappa(E):=e^{-c_1(E)/r}\operatorname{ch}(E)
 =t+e^{-a_dH_d}P_d(H_d).
 \tag{QC2.6}
\]

All terms other than \(t\) are polarization powers. The coefficients
\(c_k\) of \(e^{-a_d z}P_d(z)\), through degree six, are

\[
\begin{array}{c|l}
0&8d\\
1&0\\
2&-(d+1)^2/(4d)\\
3&(d-1)(d+1)^2/(24d^2)\\
4&-(d-3)(d+1)^2(3d-1)/(768d^3)\\
5&(d-1)(d+1)^2(d^2-6d+1)/(3840d^4)\\
6&-(d+1)^2(d^2-10d+5)(5d^2-10d+1)/(368640d^5).
\end{array}
\tag{QC2.7}
\]

\subsection{SF-QC03. A complete explicit regularity
proof}\label{SF-QC03}

The sheaf \(E\) is 112-regular with respect to \(H_b\). More strongly,
for any integer \(s\ge 1\) and any line bundle \(M\) on A satisfying

\[
 c_1(M)-112s^2H_b\quad\hbox{is nef},
 \tag{QC3.1}
\]

one has \(H^i(A,[s]^*E\otimes M)=0\) for all \(i>0\). Arbitrary flat
factors of \(M\) are allowed.

The regularity convention is \(H^i(A,F((q-i)H_b))=0\) for every \(i>0\).
The standard Castelnuovo--Mumford lemma is applied to the pushforward
under the embedding defined by \(H_b\): a zero-regular sheaf is globally
generated, remains regular under nonnegative twists, and its section
multiplication maps are surjective. An inspected published locator is
Pareschi--Popa, M-regularity and the Fourier--Mukai transform, Theorem
1.2 (author version p.2); the publisher version is Pure Appl. Math.
Q.4(3),587--611 (2008). The relevant Buchsbaum--Rim terms are in
Raicu--Sam, Compositio158(2022), pp.41--42; their exactness in the
locally split case used here is also proved below.

\textbf{Step 1: the dual bundles.} For every integer \(q\ge 1\),

\[
 H^i(A,W^*(qH_b))=H^i(A,U^*(qH_b))=0\quad(i>0).
 \tag{QC3.2}
\]

For \(W\), relative duality identifies \(W^*\) with \(R^3\pi_*L^{-1}\).
Leray and Serre duality on the ninefold \(Y\) identify the dual of
\(H^i(A,W^*(qH_b))\) with \(H^{6-i}(Y,\mathcal L(-qH_b))\). Project
\(Y\) onto \(J_y\). On an A-fibre the line bundle has matrix
\(\operatorname{diag}(1-3q,-3q)\), hence is antiample. Its cohomology is
concentrated in degree six, so total cohomology below six vanishes.

For a summand of \(U\) supported on \(C\times A\), relative duality
first introduces the relative canonical bundle; Serre duality on the
sevenfold cancels it and again gives
\(H^{6-i}(C\times A,\mathcal L_C(-qH_b))\). Project onto \(C\). The
first curve type has the same fibre matrix. For the second type use
\(z=y+x\) as the fixed curve coordinate; its matrix is

\[
 \begin{pmatrix}1-3q&1\\1&-3q\end{pmatrix}.
\]

The largest eigenvalue of the matrix before subtracting 3qI is
\((1+\sqrt 5)/2<3\), proving antiampleness for \(q\ge 1\). Flat twists
and translations do not change this calculation. Thus \(W^*,U^*\) are
7-regular.

\textbf{Step 2: a kernel estimate.} Let \(e=h^{0}(U^*(7H_b))\).
Evaluation gives

\[
 0\to K_U\to O_A(-7H_b)^e\to U^*\to0.
\]

The kernel is a vector bundle and is 14-regular. For \(i=2,\ldots ,6\),
twist by \((14-i)H_b\) and use (QC3.2) and the positivity of
\((7-i)H_b\). For \(i=1\), section multiplication from 7-regularity
gives surjectivity at twist \(13H_b\); the multiplication by degree-six
ambient polynomials factors through \(H^{0}(A,6H_b)\), which is enough
for this evaluation map.

Compose evaluation with \(U^*\to W^*\), and let \(J_{0}\) be the
resulting kernel. There is an exact sequence

\[
 0\to K_U\to J_0\to Q^*\to0.
 \tag{QC3.3}
\]

The Buchsbaum--Rim complex resolving \(J_{0}\) has terms, beginning with
a surjection \(C_{1}\to J_{0}\),

\[
 C_k=\bigwedge^{8+k}(O_A(-7H_b)^e)
       \otimes\operatorname{Sym}^{k-1}W\otimes\det W,
 \qquad 1\le k\le e-8.
 \tag{QC3.4}
\]

Its determinant factor is required globally. To verify exactness in this
case, locally split the surjection as \(K\oplus V\to V\), with dim
\(V=8\). Splitting each wedge by the number a of factors from \(K\)
identifies the residual strands, after the determinant pairing, with the
polynomial de Rham complex of total degree \(a-1\) on \(V^*\). For
\(a\ge 2\) the Euler contraction divides by the nonzero integer \(a-1\)
and gives a contracting homotopy. The \(a=1\) strand maps isomorphically
to \(K\), and the remaining \(V\) maps isomorphically to the target.
This proves the needed local exactness in characteristic zero.

The theorem of the square gives
\(W=O_A(\Theta _x)\otimes b^*\operatorname{FM}(O_J(2\Theta ))\), where
\(b(x,\alpha )=x-\alpha\). Pulling the Poincaré family back by the
polarization isogeny {[}2{]} identifies it with
\(a^*O(2\Theta )\otimes p_{2}^*O(-2\Theta )\). Taking direct image gives
\(H^{0}(O(2\Theta ))\otimes O(-2\Theta )\). Consequently

\[
 [2]^*W\simeq Q_0^{\oplus8},\qquad
 c_1(Q_0)=4\Theta_x-2\Theta_{x-\alpha}.
 \tag{QC3.5}
\]

A different symmetric normalization adds a flat factor and affects no
vanishing assertion. The matrix of \(Q_{0}\) is
\(\bigl(\begin{smallmatrix}2&2\\2&-2\end{smallmatrix}\bigr)\), with
eigenvalues \(\pm 2\sqrt 2\), while \([2]^*H_b\) has matrix 12I. Thus
\(q[2]^*H_b+(k+7)c_{1}(Q_{0})\) is ample for \(q\ge k+7\). Kodaira
vanishing on A, and trace splitting for the finite étale map {[}2{]},
prove

\[
 \operatorname{reg}_{H_b}(\operatorname{Sym}^{k-1}W\otimes\det W)
 \le k+13.
\]

Hence reg \(C_k\le 7(8+k)+k+13=69+8k\). Only the first six terms of
(QC3.4) are needed. Successive cohomology sequences show the vanishings
required for R-regularity of \(J_{0}\) whenever

\[
 R\ge\max_{1\le k\le6}(\operatorname{reg} C_k-k+1)
 \le\max_{1\le k\le6}(70+7k)=112.
\]

The final possible cohomology is in degree seven and vanishes on the
sixfold. Sequence (QC3.3), together with 14-regularity of \(K_U\), gives
112-regularity of \(Q^*\). This bound is independent of the number \(N\)
of curves and of the number \(e\) of evaluation sections.

\textbf{Step 3: the torsion dual.} On one component
\(C\times C\times \widehat{J}\), put \(F_{1}={p}\times C\),
\(F_{2}=C\times {p}\), \(S=F_{1}+F_{2}\), and let \(\Delta _C\) be the
diagonal. The sum map \(C\times C\to J\) pulls \(\Theta\) back to

\[
 A_C=4S-\Delta_C.
\]

Indeed \(\Theta |C\) has degree three and its mixed Poincaré class is
\(S-\Delta _C\). Restriction to the diagonal gives degree twelve,
agreeing with \([2]^*\Theta |C\).

The divisor \(R_C=3S-\Delta _C\) is nef. Projection of the plane quartic
from a point of \(C\) gives a degree-three morphism to \(P^{1}\). The
residual to the diagonal in its fibre correspondence is effective of
class \(R_C\). If its monodromy is \(S_{3}\), this residual divisor is
irreducible, and its self-intersection is \(18-12-4=2\). An effective
irreducible divisor of nonnegative square on a smooth surface is nef. If
the transitive monodromy is \(A_{3}\), the cover is cyclic, and the
residual divisor is the sum of the two nonidentity automorphism graphs.
Each graph has square \(-4\). Riemann--Hurwitz gives five fixed points;
the two graphs meet transversely there because the two nontrivial
eigenvalues of the order-three local action differ. Their sum has
intersection one with each graph and nonnegative intersection with any
other irreducible curve. It is again nef.

Numerically on this component,

\[
 c_1(\mathcal L_T)=3F_1+3F_2-\mathcal P_{c_1,\alpha},\qquad
 f^*H_b=3A_C+3\Theta_\alpha.
\]

For \(q\ge 3\) decompose

\[
\begin{split}
 qf^*H_b-c_1(\mathcal L_T)
 ={}&[3qA_C-6F_1-3F_2]+(3q-1)\Theta_\alpha\\
 &+[3F_1+\Theta_\alpha+\mathcal P_{c_1,\alpha}].
\end{split}
\]

The last bracket is a pullback of \(\Theta\), hence nef. The first is
\(3qR_C+(3q-6)F_{1}+(3q-3)F_{2}\), ample on \(C\times C\). Together with
the strictly positive theta class on \(\widehat{J}\) it is ample on the
product. Thus the entire class is ample. Kodaira vanishing on the smooth
fivefold gives
\(H^i(T,\omega_T\otimes\mathcal L_T^{-1}\otimes f^*O(qH_b))=0\) for
\(i>0,q\ge 3\). Therefore \(T^D\) is 8-regular. Extension (QC2.3) proves
the asserted 112-regularity of \(E\).

\textbf{Step 4: the stronger bound (QC3.1).} Pull all exact complexes
back by {[}s{]} and tensor with \(M\). In Step 1 the two output fibre
matrices are multiplied by \(s^{2}\); subtracting \(c_{1}(M)\) makes
them negative definite under (QC3.1). Thus \([s]^*U^*\otimes M\) and
\([s]^*W^*\otimes M\) have no higher cohomology. In the pulled-back
evaluation sequence the middle term is a sum of ample line bundles
\(M(-7s^{2}H_b)\), so \(H^j([s]^*K_U\otimes M)=0\) for \(j\ge 2\). No
claim about \(H^{1}\) of this kernel is needed.

For \(k\le 6\) pull \([s]^*C_k\otimes M\) back once more by {[}2{]}. Its
summands have class

\[
 [2]^*c_1(M)-7s^2(8+k)[2]^*H_b+s^2(k+7)c_1(Q_0).
\]

After removing a nef class the matrix is \(s^{2}\) times

\[
 12[112-7(8+k)]I+(k+7)\begin{pmatrix}2&2\\2&-2\end{pmatrix}.
\]

Its least eigenvalue is positive for \(1\le k\le 6\); the worst
displayed lower bound is \(168-26\sqrt 2>0\). Kodaira vanishing, trace
splitting, the six-term truncation and (QC3.3) give the required
vanishings for \([s]^*Q^*\otimes M\).

For \(T^D\), form the finite étale base change \(T'=T\times_{A,[s]}A\)
with \(g:T'\to T\) and \(f':T'\to A\). The relevant line class is

\[
 f'^*c_1(M)-g^*c_1(\mathcal L_T)
 =g^*(3f^*H_b-c_1(\mathcal L_T))
       +f'^*(c_1(M)-3s^2H_b).
\]

The first summand is ample by Step 3 and the second is nef. Kodaira
vanishing on \(T'\), finite base change, and (QC2.3) prove (QC3.1) for
\(E\).

\subsection{SF-QC04. Ordinary normalization and coherent
equivariance}\label{SF-QC04}

Let \(r=8d\) and define the coherent sheaf

\[
 B_0=[r]^*E\otimes D^{-r},\qquad D=\det E.
 \tag{QC4.1}
\]

It has rank \(r\), projective dimension at most one, \(c_{1}(B_{0})=0\),
and complete character

\[
 \operatorname{ch}(B_0)=[r]^*\kappa(E)
       =r^6t+\sum_{k=0}^6 r^{2k}c_k H_d^k.
 \tag{QC4.2}
\]

Here \(r^{6}\) is the degree-six cohomological scaling, not the isogeny
degree. Since \(\operatorname{gcd}(r,N)=1\), this normalization
preserves a strong equivariant semiregularity structure for an actual
finite subgroup of \(A\times \widehat{A}\), as follows.

Write the original group action as \(T_g=t_{a_g}^*\otimes P_g\), with
\(a_g\in A[N]\), \(P_g\in \widehat{A}[N]\), and all associativity
constraints transported from descent. Let \(a'_g\) be the unique
N-torsion point with \(r\) \(a'_g=a_g\). Define
\(T'_g=t_{a'_g}^*\otimes[r]^*P_g\). Its coherent composition maps are
the pullbacks of those for \(T_g\), since
\(t_{a'_g}^*[r]^*=[r]^*t_{a_g}^*\). The deck group \(A[r]\) acts on
\([r]^*E\) in the canonical way; each pulled-back \(P_g\) has a
canonical deck linearization. Thus these actions commute coherently and
give an action of \(A[r]\times G\). The map to \(A\times \widehat{A}\)
is

\[
 (b,g)\longmapsto(b+a'_g,\,[r]^*P_g).
\]

It is injective: the first summands have coprime r- and N-torsion
orders, and multiplication by \(r\) is injective on N-torsion in the
dual. There is no invisible central kernel. Descent for the deck group
identifies the invariant \(\operatorname{Ext}^{2}\) with
\(\operatorname{Ext}^{2}(E,E)\), and then taking G-invariants gives the
original semiregular invariant space. Trace pullback is injective under
the isogeny. Tensor by \(D^{-r}\) conjugates this honest action by an
autoequivalence, hence preserves coherence and injectivity of its
underlying subgroup. This proves strong, and therefore weak, equivariant
semiregularity of \(B_{0}\). Ordinary negative Ext vanishes because
\(B_{0}\) is a coherent sheaf.

The full character (QC4.2) stays Hodge on the split Weil period family
with polarization \(h=3H_d\). On a suitable finite level cover the
rational Weil plane is constant, and \(r^{6}t\) extends as a flat Hodge
class. Apply Perry Theorem 1.2 with B-field zero. Its proof gives an
ordinary perfect deformation, not merely a twisted one: in §6.2 choose
\(b_{0}=0,n=1\) and the zero gerbe class. Corollary 5.17 may replace its
band by \(\mu _q\), but pushing out the zero class keeps the Brauer
class zero. Theorem 6.3 returns the object in the ambient category
\(\operatorname{Perf}(A_U,0)\). Any gerbe on an auxiliary quotient used
to geometrize the invariant category does not change this ambient
conclusion.

After shrinking the étale neighbourhood, the ordinary perfect
deformation is a flat coherent sheaf B and has fibrewise projective
dimension at most one. Here are precise algebraic details of the
openness step. At every point of the special fibre the derived
residue-field complex has amplitude in \([-1,0]\), because \(B_{0}\) has
projective dimension at most one. Openness of the Betti-number
vanishings for a perfect complex (Stacks Lemma36.31.1, Tag0BDI)
therefore gives a two-term vector-bundle presentation in degrees
\(-1,0\) on a neighbourhood of that point. On the special fibre its
differential is injective as a map of sheaves, since \(B_{0}\) is a
sheaf. The local criterion in Stacks Lemma10.99.1, Tag00ME, applied to
the local rings of the base and total space, proves that the
differential is injective near this point and that its cokernel is flat
over the base there. Openness of the flat locus (Theorem37.15.1,
Tag0399, citing EGA IV\(_3\), 11.3.1) extends this to an open
neighbourhood of the entire special fibre. Properness permits shrinking
the base so that every point of the new total space lies in this
neighbourhood. The two-term presentation gives \(pd\le 1\) on every
fibre, and the cokernel gives the desired ordinary flat coherent family.
No global boundedness of all sheaves with this Chern character is
asserted.

Put

\[
 R=224r^2+r(d+2)+6=14344d^2+16d+6.
 \tag{QC4.3}
\]

Then \(B_{0}\) is R-regular with respect to \(h\). Indeed
\(2h-H_b=3((2d-1)x+y)\) is ample. The matrix of \((d+2)h-c_{1}(D)\) is

\[
 \begin{pmatrix}d(d+8)&4d\\4d&d+8\end{pmatrix},
\]

whose determinant is \(d((d+8)^{2}-16d)=d(d^{2}+64)>0\). For
\(i=1,\ldots ,6\), the line bundle \(M_i=D^{-r}\otimes h^{R-i}\)
satisfies

\[
 c_1(M_i)-112r^2H_b
 =112r^2(2h-H_b)+r((d+2)h-c_1(D))+(6-i)h,
\]

which is ample. SF-QC03 supplies all six vanishings. Upper
semicontinuity in the ordinary flat coherent family preserves them after
shrinking. Thus every \(B_s\) on this open is R-regular, of rank \(r\),
\(c_{1}=0\) and \(pd\le 1\), with the same full character.

\subsection{SF-QC05. A numerical effective
detector}\label{SF-QC05}

Define the explicit integers

\[
 \begin{split}
 r&=8d,\qquad R=14344d^2+16d+6,\\
 V&=\int_A h^6=524880d^3,\\
 M&=729d^3r(R+r^2)^6,\\
 A_2&=\frac{1024}{9}d^3(d+1)^2.
 \end{split}
 \tag{QC5.1}
\]

\(A_{2}\) is an integer because \(d=4mL^{2}\) is divisible by nine. Set

\[
 D_{\rm CM}(m)=V\left((r+2)R A_2
       +7M\binom{r+1}{2}R^2+343\binom M3\right).
 \tag{QC5.2}
\]

This is a numerical formula involving only \(m\), powers, binomial
coefficients and the explicit finite lcm defining \(L\). On the entire
fixed split Weil family containing (A,h,K), there is an effective
integral three-cycle with nonzero Weil projection and \(h\)-degree at
most \(D_{\rm CM}(m)\), using the authorized detector closedness input
to pass from the open obtained above to its component. No Chow
denominator occurs in the detecting cycle.

In particular the simpler safe integer

\[
 \boxed{D_{\rm CM}(m)\le 2^{318}d^{51}
                  =2^{420}m^{51}L^{102}}
 \tag{QC5.2a}
\]

may be used instead. Indeed \(r+2\le 2^{4}d\), \(R<2^{14}d^{2}\),
\(A_{2}<2^{9}d^{5}\), \(V<2^{20}d^{3}\) and
\(\operatorname{binom}(r+1,2)\le 2^{6}d^{2}\). Also
\(R+r^{2}\le 14430d^{2}<2^{14}d^{2}\), so \(M<2^{97}d^{16}\). The three
summands in (QC5.2) are respectively bounded by \(2^{47}d^{11}\),
\(2^{154}d^{25}\), and \(2^{317}d^{51}\); for the last use 343
\(\operatorname{binom}(M,3)<64M^{3}\). Since \(d\ge 1\), their sum is
below \(2^{318}d^{51}\).

\textbf{Proof.} On the deformation open, \(F_s=B_s(Rh)\) is zero-regular
and globally generated. Its higher cohomology vanishes and
\(v=h^{0}(F_s)=\chi (F_s)\) is constant. The evaluation sequence

\[
 0\to K_s\to O_{A_s}^{v}\to F_s\to0
\]

has a vector-bundle kernel because \(\operatorname{pd}(F_s)\le 1\). Put
\(k=v-r\). The same section multiplication argument used above shows
that \(K_s\) is 7-regular: for cohomological degrees at least two use
the positivity of \((7-i)h\) and zero-regularity of \(F_s\), and in
degree one use surjectivity of
\(H^{0}(F_s)\otimes H^{0}(O(6h))\to H^{0}(F_s(6h))\). Hence
\(G_s=K_s(7h)\) is a globally generated vector bundle.

Let

\[
 A_3=\frac{65536}{81}d^4(d-1)(d+1)^2.
\]

From (QC2.7) and \(c_{1}(B_s)=0\),

\[
 c_2(B_s)=A_2h^2,\qquad
 c_3(B_s)=A_3h^3+2r^6t_s.
 \tag{QC5.3}
\]

The sign and all coefficients here follow from
\(\operatorname{ch}_{2}=-c_{2}\) and \(\operatorname{ch}_{3}=c_{3}/2\)
when \(c_{1}=0\). Inverting \(c(F_s)\) in the evaluation sequence and
then twisting by 7h gives

\[
\begin{split}
 c_3(G_s)= {}&-2r^6t_s+ S_kh^3,\\
 S_k={}&-A_3+(r+2)R A_2-\binom{r+2}{3}R^3\\
 &+7(k-2)\left(\binom{r+1}{2}R^2-A_2\right)
   -49rR\binom{k-1}{2}+343\binom k3.
\end{split}
\tag{QC5.4}
\]

Thus \(c_{3}(G_s)\) has nonzero Weil projection. In particular
\(k\ge 3\). Choose \(k-2\) general global sections of \(G_s\). Their
first rank-drop locus has expected codimension three, or is empty,
because the bundle is globally generated and the universal evaluation is
transverse to the matrix-rank strata. The Thom--Porteous formula gives
its class \(c_{3}(G_s)\). Nonzero Weil projection excludes the empty
case. It is therefore an integral effective detecting cycle.

It remains to bound \(k\). Write \(c_j\) for (QC2.7). Since \(H_d t=0\)
and \(td(A)=1\), Riemann--Roch gives

\[
 v=V\sum_{j=0}^6\frac{r^{2j}c_j}{3^j}
                      \frac{R^{6-j}}{(6-j)!}.
\]

For \(d\ge 1\) each coefficient of \(P_d\) has absolute value at most
r/j!. Since \(0\le a_d<1/4\), multiplying by
\(\operatorname{exp}(-a_dz)\) gives \(|c_j|\le r(5/4)^j/j!\).
Consequently

\[
 v\le\frac{Vr}{6!}(R+5r^2/12)^6
    \le729d^3r(R+r^2)^6=M.
\]

Thus \(3\le k\le M\). Since \(d>1\), both \(A_{2},A_{3}\) are positive.
Discarding the negative terms in (QC5.4) and replacing \(k\) by \(M\) in
the remaining positive terms gives \(S_k\le (r+2)RA_{2}+7M\)
\(\operatorname{binom}(r+1,2)R^{2}+343\) \(\operatorname{binom}(M,3)\).
Finally \(h^{3}t_s=0\), so the degree is \(VS_k\le D_{\rm CM}(m)\),
proving (QC5.2).

The source polarization is specified: \(h=3(dx+y)\) has type
(3,3,3,3d,3d,3d), product \(729d^{3}\). This result must not be
described as a principal-polarization bound without an additional
lattice/isogeny calculation. For any fixed target split polarized
K-lattice, a chosen rational polarization isometry whose denominators
are cleared by \(q\) gives a fixed isogeny with pullback target
polarization \(q^{2}h\); pushforward transports the detector bound with
cost \(q^{6}\). A numerical bound uniform over all target polarization
products requires bounding that \(q\), separately.

For class normalization, the exterior pairing gives
\(\int t^{2}=-8dN^{2}\): each symplectic plane contributes
\(Q_uQ_{-u}=-4d\) times its oriented volume, and
\(2(ab)^{2}=2N^{2}/(16d^{2})\). Therefore, for \(t_B=r^{6}t\),

\[
 -\int t_B^2=8dN^2r^{12}=2^{39}d^{13}N^2.
\]

Moreover \(A_{3}\) is an integer for the chosen \(d\), since 9 divides
\(d\). Formula (QC5.3) then expresses \(2t_B\) as the difference of
integral algebraic cycle classes. Thus the lattice denominator
\(\delta =2\) is valid for this explicit flat Weil class. This does not
assert that an arbitrary rational Weil vector has denominator two.

\subsection{Source and status notes}\label{source-and-status-notes}

\begin{itemize}
\item
  The finite torsion avoidance and the denominator-free effective Chern
  construction are new arguments in this run. The proof does not invoke
  an unverified generic-rank or generic-position statement.
\item
  Pareschi--Popa, M-regularity and the Fourier--Mukai transform, Pure
  Appl. Math. Q.4(3),587--611 (2008), Theorem1.2; exact basic regularity
  statement inspected in the author PDF p.2. Publisher
  DOI10.4310/PAMQ.2008.v4.n3.a1.
\item
  Raicu--Sam, Bi-graded Koszul modules, K3 carpets, and Green's
  conjecture, Compositio158(2022),33--56, pp.41--42 immediately before
  Theorem3.2. The globally correct Buchsbaum--Rim determinant factor was
  inspected in the published PDF, and the locally split exactness is
  proved here.
\item
  Bauer--Szemberg, On tensor products of ample line bundles on abelian
  varieties, Math.Z.223(1996),79--85, Theorem1.1(a),(b), author version
  p.1. The proof and publication header were inspected.
\item
  Fulton, Intersection Theory, second edition (1998), Theorem14.4,
  pp.254--255. The Thom--Porteous statement used here was checked
  against the publisher's chapter abstract and the indexed primary book
  text; a full independent audit of the entire chapter proof is not
  claimed.
\item
  Mukai1981, Theorem2.2 and equation3.1, and Buchweitz--Flenner2003
  Proposition4.2, remain the published Fourier--Mukai and trace inputs
  already checked in D-SR-01--04.
\item
  Perry2604.00511v2 remains an authorized preprint input. The zero-B
  ordinary conclusion is an inference from the full proofs of
  Corollary5.17, Theorems6.1/6.3 and Theorem1.2 in §6.2, not an extra
  theorem quoted without checking.
\item
  Exact Perry page locators for this inference: Corollary5.17 pp.20--21;
  Theorems6.1/6.3 pp.25--26; §6.2 and the proof of Theorem1.2 pp.26--28.
  All were inspected in the official versioned PDF.
\item
  Stacks Project Tags0BDI,00ME,0399 and07VK were opened and checked for,
  respectively, openness of perfect-complex Betti vanishings, the
  injectivity/flat-cokernel criterion, openness of the flat locus, and
  perfect cohomology with arbitrary base change. The last two provide
  the precise family hypotheses used for upper semicontinuity.
\item
  All regularity constants used for the numerical estimate have been
  proved above. HECKE\_DIV M7 is provenance for the architecture of that
  argument, not a premise, and Markman's output sheaf and
  general-position assertions are unnecessary here.
\end{itemize}
